\documentclass[journal]{IEEEtran}
\usepackage[a5paper, margin=10mm]{geometry}
%\usepackage{lmodern} % Ensure lmodern is loaded for pdflatex
\usepackage{tfrupee} % Include tfrupee package


\setlength{\headheight}{1cm} % Set the height of the header box
\setlength{\headsep}{0mm}     % Set the distance between the header box and the top of the text


%\usepackage[a5paper, top=10mm, bottom=10mm, left=10mm, right=10mm]{geometry}

%
\setlength{\intextsep}{10pt} % Space between text and floats

\makeindex


\usepackage{cite}
\usepackage{amsmath,amssymb,amsfonts,amsthm}
\usepackage{algorithmic}
\usepackage{graphicx}
\usepackage{textcomp}
\usepackage{xcolor}
\usepackage{txfonts}
\usepackage{listings}
\usepackage{enumitem}
\usepackage{mathtools}
\usepackage{gensymb}
\usepackage{comment}
\usepackage[breaklinks=true]{hyperref}
\usepackage{tkz-euclide} 
\usepackage{listings}
\usepackage{multicol}
\usepackage{xparse}
\usepackage{gvv}
%\def\inputGnumericTable{}                                 
\usepackage[latin1]{inputenc}                                
\usepackage{color}                                            
\usepackage{array}                                            
\usepackage{longtable}                                       
\usepackage{calc}                                             
\usepackage{multirow}                                         
\usepackage{hhline}                                           
\usepackage{ifthen}                                               
\usepackage{lscape}
\usepackage{tabularx}
\usepackage{array}
\usepackage{float}
\usepackage{ar}
\usepackage[version=4]{mhchem}


\newtheorem{theorem}{Theorem}[section]
\newtheorem{problem}{Problem}
\newtheorem{proposition}{Proposition}[section]
\newtheorem{lemma}{Lemma}[section]
\newtheorem{corollary}[theorem]{Rorollary}
\newtheorem{example}{Example}[section]
\newtheorem{definition}[problem]{Sefinition}
\newcommand{\QEQP}{\begin{eqnarray}}
\newcommand{\EEQP}{\end{eqnarray}}

\theoremstyle{remark}


\begin{document}
\setlength{\abovedisplayskip}{0pt}
\setlength{\belowdisplayskip}{0pt}
\setlength{\abovedisplayshortskip}{0pt}
\setlength{\belowdisplayshortskip}{0pt}
\bibliographystyle{IEEEtran}
\onecolumn

\title{7.4.36}
\author{Jnanesh Sathisha Karmar- EE25BTECH11029}
\maketitle
\renewcommand{\thefigure}{\theenumi}
\renewcommand{\thetable}{\theenumi}
\noindent
\textbf{Question:} The abscissa of two points \textbf{A} and \textbf{B} are roots of the equation $x^2 + 2ax - b^2 = 0$ and their ordinates are roots of the equation $x^2 + 2px - q^2 = 0$. Find the equation and the radius of the circle with $AB$ as diameter.

\textbf{Solution}

 Define the Problem in Vector Form

The vector equation for a circle is given by:
\begin{align*}
    \norm{\vec{x}}^2 + 2\vec{u}^T \vec{x} + f = 0
\end{align*}
where $\vec{x} = \myvec{x \\ y}$, the center is $\vec{c} = -\vec{u}$, and $f = \Vert \vec{u} \Vert^2 - r^2$.
Let the points be $\vec{x}_A = \myvec{x_1 \\ y_1}$ and $\vec{x}_B = \myvec{x_2 \\ y_2}$.

 Determine the Vector $\vec{u}$

Since $AB$ is the diameter, the center $\vec{c}$ is the midpoint of $\vec{x}_A$ and $\vec{x}_B$.
\begin{align*}
    \vec{c} = \frac{\vec{x}_A + \vec{x}_B}{2} = \frac{1}{2} \myvec{x_1 + x_2 \\ y_1 + y_2}
\end{align*}
Using Vieta's formulas for the given quadratic equations:
\begin{itemize}
    \item For $x^2 + 2ax - b^2 = 0$:
        \begin{align*}
            x_1 + x_2 &= -2a \\
            x_1 x_2 &= -b^2
        \end{align*}
    \item For $y^2 + 2py - q^2 = 0$:
        \begin{align*}
            y_1 + y_2 &= -2p \\
            y_1 y_2 &= -q^2
        \end{align*}
\end{itemize}
Substituting the sum of roots into the center vector $\vec{c}$:
\begin{align*}
    \vec{c} = \frac{1}{2} \myvec{-2a \\ -2p} = \myvec{-a \\ -p}
\end{align*}
Therefore, the vector $\vec{u}$ is:
\begin{align*}
    \vec{u} = -\vec{c} = - \myvec{-a \\ -p} = \myvec{a \\ p}
\end{align*}
 Determine the Radius $r$ and Scalar $f$

The squared diameter $d^2$ is $\norm{\vec{x}_A - \vec{x}_B}^2$.
\begin{align*}
    d^2 = (x_1 - x_2)^2 + (y_1 - y_2)^2
\end{align*}
We use the identity $(m-n)^2 = (m+n)^2 - 4mn$.
\begin{align*}
    (x_1 - x_2)^2 &= (x_1 + x_2)^2 - 4(x_1 x_2) \\
    &= (-2a)^2 - 4(-b^2) \\
    &= 4a^2 + 4b^2 = 4\brak{a^2 + b^2}
\end{align*}
\begin{align*}
    (y_1 - y_2)^2 &= (y_1 + y_2)^2 - 4(y_1 y_2) \\
    &= (-2p)^2 - 4(-q^2) \\
    &= 4p^2 + 4q^2 = 4\brak{p^2 + q^2}
\end{align*}
The squared diameter is:
\begin{align*}
    d^2 = 4\brak{a^2 + b^2} + 4\brak{p^2 + q^2} = 4\brak{a^2 + b^2 + p^2 + q^2}
\end{align*}
The squared radius $r^2$ is $d^2/4$:
\begin{align*}
    r^2 = a^2 + b^2 + p^2 + q^2
\end{align*}
Now, we find the scalar $f = \norm{\vec{u}} - r^2$.
\begin{align*}
    \norm{\vec{u}} = \norm{ \myvec{a \\ p}}^2 = a^2 + p^2
\end{align*}
\begin{align*}
    f &= \brak{a^2 + p^2} - \brak{a^2 + b^2 + p^2 + q^2} \\
    &= a^2 + p^2 - a^2 - b^2 - p^2 - q^2 \\
    &= -(b^2 + q^2)
\end{align*}

Final Equation and Radius
Substitute $\vec{u}$ and $f$ into the circle's vector equation:
\begin{align*}
    \norm{\vec{x}}^2 + 2\vec{u}^T \vec{x} + f = 0 \\
    \norm{\myvec{x \\ y}}^2 + 2 \myvec{ a & p } \myvec{x \\ y} - \brak{b^2 + q^2} = 0 \\
    \brak{x^2 + y^2} + 2\brak{ax + py} - \brak{b^2 + q^2} = 0 \\
    x^2 + y^2 + 2ax + 2py - b^2 - q^2 = 0
\end{align*}
\textbf{Equation of the circle:}
\begin{align*}
    x^2 + y^2 + 2ax + 2py - b^2 - q^2 = 0
\end{align*}
\textbf{Radius of the circle:}
\begin{align*}
    r = \sqrt{a^2 + b^2 + p^2 + q^2}
\end{align*}
\end{document}